% Compile with XeLaTeX (Overleaf: Menu > Compiler > XeLaTeX).
% Logo: upload logo-sgu.png cùng thư mục với file .tex trên Overleaf.
% Burndown: tự vẽ bằng TikZ/pgfplots, không cần file ảnh ngoài.
\documentclass[12pt,a4paper]{report}
\usepackage{fontspec}
\setmainfont{texgyretermes-regular.otf}[BoldFont=texgyretermes-bold.otf,ItalicFont=texgyretermes-italic.otf,BoldItalicFont=texgyretermes-bolditalic.otf]
\setsansfont{texgyreheros-regular.otf}[BoldFont=texgyreheros-bold.otf,ItalicFont=texgyreheros-italic.otf,BoldItalicFont=texgyreheros-bolditalic.otf]
\setmonofont{lmmono10-regular.otf}
\usepackage[top=2.3cm,bottom=2.3cm,left=2.7cm,right=2.2cm]{geometry}
\usepackage{graphicx,array,longtable,multirow}
\usepackage{setspace,titlesec,fancyhdr,needspace}
\usepackage{tikz}
\usepackage{pgfplots}
\pgfplotsset{compat=1.18}
\usepackage[unicode,hidelinks]{hyperref}
\usepackage{xurl}
\setstretch{1.25}
\setlength{\parindent}{0pt}
\setlength{\parskip}{6pt}
\setlength{\emergencystretch}{2em}
\renewcommand{\contentsname}{MỤC LỤC}
\renewcommand{\chaptername}{CHƯƠNG}
\setcounter{tocdepth}{1}
\setcounter{secnumdepth}{1}
\titleformat{\chapter}[block]{\bfseries\Large}{CHƯƠNG \thechapter}{0.7em}{}
\titlespacing*{\chapter}{0pt}{0pt}{16pt}
\titleformat{\section}{\bfseries\large}{\thesection}{0.6em}{}
\titlespacing*{\section}{0pt}{12pt}{7pt}
\pagestyle{fancy}
\fancyhf{}
\fancyhead[L]{\small BrewLite\quad |\quad Báo cáo Assignment 1}
\fancyfoot[C]{Trang \thepage}
\renewcommand{\headrulewidth}{0pt}
\fancypagestyle{plain}{\fancyhf{}\fancyfoot[C]{Trang \thepage}\renewcommand{\headrulewidth}{0pt}}
\begin{document}
\begin{titlepage}
\centering
{\bfseries\large TRƯỜNG ĐẠI HỌC SÀI GÒN\\KHOA CÔNG NGHỆ THÔNG TIN\par}
\vspace{0.6cm}
\includegraphics[width=2.6cm]{logo-sgu.png}
\vspace{0.7cm}

{\bfseries\Large BÁO CÁO ĐỒ ÁN MÔN CÔNG NGHỆ PHẦN MỀM\par}
\vspace{1cm}
{\bfseries\large BREWLITE\\HỆ THỐNG ĐẶT VÀ THANH TOÁN ĐỒ UỐNG TẠI QUẦY\par}
\vspace{1cm}
Giảng viên hướng dẫn: ThS. Từ Lãng Phiêu\\
Nhóm: 05
\vspace{0.6cm}


\begingroup
\small
\setlength{\tabcolsep}{5pt}
\renewcommand{\arraystretch}{1.2}
\begin{longtable}{|>{\centering\arraybackslash}p{\dimexpr 0.09676\linewidth-2\tabcolsep-2\arrayrulewidth\relax}|>{\raggedright\arraybackslash}p{\dimexpr 0.55908\linewidth-2\tabcolsep-2\arrayrulewidth\relax}|>{\raggedright\arraybackslash}p{\dimexpr 0.34416\linewidth-2\tabcolsep-2\arrayrulewidth\relax}|}
\hline
\textbf{STT} & \textbf{Họ và tên thành viên} & \textbf{Mã số sinh viên (MSSV)} \\ \hline
\endfirsthead
\hline
\textbf{STT} & \textbf{Họ và tên thành viên} & \textbf{Mã số sinh viên (MSSV)} \\ \hline
\endhead
\endfoot
\endlastfoot
1 & Nguyễn Minh Hiển (Trưởng nhóm) & 3124410088 \\ \hline
2 & Trần Đăng Khoa & 3124410161 \\ \hline
3 & Nguyễn Chí Thành & 3124410324 \\ \hline
4 & Phạm Minh Thiện & 3124410336 \\ \hline
5 & Trần Kim Xuyến & 3124410417 \\ \hline
6 & Trần Minh Hiệp & 3124410090 \\ \hline
\end{longtable}
\endgroup


\vfill

{\bfseries THÀNH PHỐ HỒ CHÍ MINH, THÁNG 11 NĂM 2026}

\end{titlepage}

\pagenumbering{roman}

\begingroup\singlespacing\small\tableofcontents\endgroup

\clearpage

\pagenumbering{arabic}

\chapter{GIỚI THIỆU TỔNG QUAN VÀ TỔ CHỨC DỰ ÁN}

BrewLite là ứng dụng web hỗ trợ đặt và thanh toán đồ uống tại quầy. Khách hàng có thể xem thực đơn, chọn kích cỡ và topping, quản lý giỏ, đăng nhập rồi thanh toán mô phỏng. Báo cáo Assignment 1 trình bày phạm vi sản phẩm, cách tổ chức nhóm, nhật ký thực hiện 10 task, kết quả kiểm thử và kết quả trình diễn, bàn giao sản phẩm.


\Needspace{8\baselineskip}
\section{Phạm vi dự án và công nghệ sử dụng}

\textbf{Phạm vi nghiệp vụ cốt lõi (MVP):} Hệ thống hỗ trợ luồng trải nghiệm khép kín dành cho khách hàng: Tra cứu thực đơn (Menu) \ensuremath{\rightarrow} Xem chi tiết \& tùy chỉnh sản phẩm \ensuremath{\rightarrow} Quản lý giỏ hàng (Cart) \ensuremath{\rightarrow} Đăng ký / Đăng nhập tài khoản \ensuremath{\rightarrow} Thực hiện Checkout \& Thanh toán trực tuyến (mô phỏng giao dịch qua Ví điện tử / Thẻ ngân hàng) \ensuremath{\rightarrow} Xem màn hình xác nhận đơn hàng và tra cứu Lịch sử mua hàng. Bên cạnh đó, dự án mở rộng thêm mô-đun quản lý dành cho Nhân viên quầy (Staff) để tiếp nhận đơn, cập nhật trạng thái chế biến và quản lý kho nguyên liệu thực tế.


\textbf{Kiến trúc công nghệ lựa chọn:}


\textbf{Frontend Framework: }Next.js (React), TypeScript, Tailwind CSS (thiết kế giao diện phản hồi), Zustand (quản lý trạng thái toàn cục cho giỏ hàng và phiên xác thực).


\textbf{Backend Framework: }NestJS (Node.js), Prisma ORM (tương tác cơ sở dữ liệu), PostgreSQL (hệ quản trị cơ sở dữ liệu quan hệ).


\textbf{Bảo mật \& Kiểm soát: }Mã hóa mật khẩu bằng Bcrypt, bảo vệ API bằng JWT Guard, chuẩn hóa dữ liệu đầu vào bằng Class-Validator / DTO.


\textbf{Đóng gói \& Triển khai: }Docker \& Docker Compose đa dịch vụ (Web, API, PostgreSQL Database) giúp đồng nhất môi trường phát triển và vận hành.


\Needspace{8\baselineskip}
\section{Phân công nhiệm vụ}

\begingroup
\small
\setlength{\tabcolsep}{5pt}
\renewcommand{\arraystretch}{1.2}
\begin{longtable}{|>{\raggedright\arraybackslash}p{\dimexpr 0.25806\linewidth-2\tabcolsep-2\arrayrulewidth\relax}|>{\raggedright\arraybackslash}p{\dimexpr 0.74194\linewidth-2\tabcolsep-2\arrayrulewidth\relax}|}
\hline
\textbf{Thành viên phụ trách} & \textbf{Phạm vi công việc và phân công chi tiết} \\ \hline
\endfirsthead
\hline
\textbf{Thành viên phụ trách} & \textbf{Phạm vi công việc và phân công chi tiết} \\ \hline
\endhead
\endfoot
\endlastfoot
\textbf{Nguyễn Minh Hiển} & Phụ trách Task T01, T06, T10B; Quản lý khởi tạo dự án, API tạo đơn hàng, cơ chế Idempotency thanh toán chống trùng lặp, trừ kho nguyên liệu và mở rộng giao diện Nhân viên. \\ \hline
\textbf{Trần Đăng Khoa} & Phụ trách Task T07, T10A; Xây dựng mô-đun Xác thực người dùng (JWT, Bcrypt), Máy trạng thái đơn hàng (State Machine), Chương trình Khuyến mãi và Điểm thưởng tích lũy. \\ \hline
\textbf{Nguyễn Chí Thành} & Phụ trách Task T02 và Backend T08; Xây dựng API Danh mục/Chi tiết Sản phẩm (NestJS/Prisma) và API Mô phỏng Thanh toán trực tuyến. \\ \hline
\textbf{Phạm Minh Thiện} & Phụ trách Task T05 và Frontend T08; Phát triển Mô-đun Giỏ hàng linh hoạt (Zustand) và Màn hình Checkout/Thanh toán phía client. \\ \hline
\textbf{Trần Kim Xuyến} & Phụ trách Task T03 và Frontend T09; Phát triển giao diện Trang thực đơn (Menu Page) và các màn hình Xác nhận đơn hàng, Lịch sử mua hàng. \\ \hline
\textbf{Trần Minh Hiệp} & Phụ trách Task T04 và DevOps T09; Xây dựng màn hình Chi tiết sản phẩm \& Tùy chỉnh topping, Đóng gói Docker Compose và tài liệu hướng dẫn bàn giao. \\ \hline
\end{longtable}
\endgroup


\Needspace{12\baselineskip}
\textbf{Tỷ lệ đóng góp của các thành viên}

Tỷ lệ đóng góp được tổng hợp cho toàn bộ quá trình thực hiện dự án, gồm phát triển chức năng, kiểm thử, sửa lỗi, tài liệu và bàn giao.

\begingroup
\small
\setlength{\tabcolsep}{6pt}
\renewcommand{\arraystretch}{1.25}
\begin{longtable}{|>{\centering\arraybackslash}p{1cm}|>{\raggedright\arraybackslash}p{8cm}|>{\centering\arraybackslash}p{3cm}|}
\hline
\textbf{STT} & \textbf{Họ và tên} & \textbf{Tỷ lệ đóng góp} \\ \hline
\endfirsthead
\hline
\textbf{STT} & \textbf{Họ và tên} & \textbf{Tỷ lệ đóng góp} \\ \hline
\endhead
\endfoot
\endlastfoot
1 & Nguyễn Minh Hiển & 18\% \\ \hline
2 & Nguyễn Chí Thành & 17\% \\ \hline
3 & Trần Đăng Khoa & 16,25\% \\ \hline
4 & Phạm Minh Thiện & 16,25\% \\ \hline
5 & Trần Kim Xuyến & 16,25\% \\ \hline
6 & Trần Minh Hiệp & 16,25\% \\ \hline
\multicolumn{2}{|c|}{\textbf{Tổng cộng}} & \textbf{100\%} \\ \hline
\end{longtable}
\endgroup

\Needspace{8\baselineskip}
\section{Tổ chức công việc theo Scrum}

Nhóm theo dõi công việc bằng GitHub Issues và Project, phát triển trên các nhánh và gửi yêu cầu review trước khi tích hợp vào nhánh chính. Kế hoạch được chia thành ba Sprint theo nhóm chức năng trong đề. Lịch và phạm vi ở bảng dưới là kế hoạch trên Project; ngày thực hiện, sửa đổi và bàn giao từng task được ghi riêng trong chương 2.


\textbf{Product Owner (PO), Scrum Master (SM): }Nguyễn Minh Hiển


\begingroup
\small
\setlength{\tabcolsep}{5pt}
\renewcommand{\arraystretch}{1.2}
\begin{longtable}{|>{\centering\arraybackslash}p{\dimexpr 0.17204\linewidth-2\tabcolsep-2\arrayrulewidth\relax}|>{\raggedright\arraybackslash}p{\dimexpr 0.25806\linewidth-2\tabcolsep-2\arrayrulewidth\relax}|>{\raggedright\arraybackslash}p{\dimexpr 0.56989\linewidth-2\tabcolsep-2\arrayrulewidth\relax}|}
\hline
\textbf{Sprint} & \textbf{Thời gian trên Project} & \textbf{Phạm vi} \\ \hline
\endfirsthead
\hline
\textbf{Sprint} & \textbf{Thời gian trên Project} & \textbf{Phạm vi} \\ \hline
\endhead
\endfoot
\endlastfoot
Sprint 1 & 14–20/10/2026 & T01–T04: nền tảng và Menu. \\ \hline
Sprint 2 & 21–27/10/2026 & T05–T07: giỏ hàng, tạo đơn, xác thực. \\ \hline
Sprint 3 & 28/10–03/11/2026 & T08–T10: thanh toán, bàn giao và nghiệp vụ nâng cao. \\ \hline
\end{longtable}
\endgroup


\Needspace{8\baselineskip}
\section{Mục tiêu của đề tài}

BrewLite hướng đến việc giúp khách chọn đồ uống và thanh toán không dùng tiền mặt ngay trên trình duyệt, giảm thao tác đặt món thủ công tại quầy. Mục tiêu của dự án là xây dựng một phiên bản có thể chạy từ đầu đến cuối, thay vì chỉ có các màn hình rời rạc.


Về kỹ thuật, nhóm sử dụng Next.js và NestJS; tổ chức dữ liệu, xác thực tài khoản và xử lý đơn hàng ở máy chủ. Về làm việc nhóm, các thành viên chia nhiệm vụ, lưu lịch sử phát triển và kiểm tra phần việc trước khi tích hợp.


\Needspace{8\baselineskip}
\section{Người dùng và phạm vi sử dụng}

Khách hàng xem Menu, chọn kích cỡ và topping, quản lý giỏ, đăng nhập, thanh toán và xem lại đơn của mình. Nhân viên quầy dùng khu vực riêng để theo dõi đơn và nguyên liệu; đây là phần mở rộng hỗ trợ vận hành, không thay thế các chức năng khách hàng thuộc yêu cầu tối thiểu.


Nhóm phát triển chịu trách nhiệm xây dựng và bàn giao bản chạy được. Giảng viên đối chiếu sản phẩm, quy trình và minh chứng với tiêu chí đề bài. Thanh toán ví/thẻ trong A1 là mô phỏng; hệ thống không thực hiện giao dịch tiền thật.


\Needspace{8\baselineskip}
\section{Yêu cầu và luồng xử lý chính}

Danh sách đồ uống phải có tên, ảnh và giá. Trang chi tiết cho chọn kích cỡ và topping, đồng thời cập nhật giá. Giỏ hàng hỗ trợ thêm, chỉnh sửa, xóa và tính tổng; các trạng thái tải, rỗng hoặc lỗi cần có thông báo để khách biết cách tiếp tục.


Khách cần đăng nhập trước khi tạo đơn. Máy chủ kiểm tra dữ liệu, tính lại số tiền và lưu đơn chờ thanh toán. Khi thanh toán mô phỏng thành công, đơn chuyển sang đã thanh toán, giỏ được xóa và trang xác nhận hiển thị mã đơn. Nếu thất bại, đơn ghi nhận lỗi thanh toán và giỏ được giữ để khách thử lại.


Lịch sử đơn chỉ trả dữ liệu thuộc tài khoản đang đăng nhập. Giá, khuyến mãi, điểm và tình trạng đơn được xác định tại máy chủ; giao diện hiển thị kết quả nhận từ API.


Yêu cầu chất lượng của đề gồm phản hồi API dưới 500 ms với dữ liệu mẫu, mật khẩu được băm, sử dụng JWT và kiểm tra đầu vào. Các kết quả đạt cần gắn với môi trường và phạm vi đã kiểm thử, không suy rộng thành bảo đảm cho mọi mức tải.


\Needspace{8\baselineskip}
\section{Tổ chức dữ liệu và xử lý nghiệp vụ}

Giao diện Next.js gửi yêu cầu đến REST API NestJS. Backend dùng Prisma để đọc, ghi PostgreSQL. Các nhóm dữ liệu chính gồm sản phẩm, tài khoản, đơn, chi tiết đơn và thanh toán; phần mở rộng có nguyên liệu, công thức, khuyến mãi và lịch sử điểm.


Đơn hàng lưu lại lựa chọn và số tiền tại thời điểm đặt. Kho được tính theo nguyên liệu cần cho món, kích cỡ và số phần topping. Khi thanh toán, hệ thống kiểm tra lại tổng nhu cầu nguyên liệu của giỏ trong giao dịch cơ sở dữ liệu; việc kiểm tra trên giao diện không có nghĩa đã giữ chỗ trong kho.


\Needspace{8\baselineskip}
\section{Quản lý tiến độ và phối hợp}

Các Issue là danh sách công việc; Sprint thể hiện nhóm nhiệm vụ theo từng giai đoạn; bản tích hợp chạy được là kết quả tăng dần của sản phẩm. Thành viên trao đổi phạm vi và phần giao nhau trước khi sửa, phát triển trên nhánh rồi gửi mã nguồn để review. Khi phát hiện lỗi, nhóm ghi nhận cách tái hiện, sửa và kiểm tra lại.


Theo Scrum, lập kế hoạch đầu Sprint dùng để chốt mục tiêu; trao đổi hằng ngày dùng để nhận biết trở ngại; Review trình diễn kết quả và Retrospective rút kinh nghiệm. Ngày 03/11/2026, nhóm tổ chức Sprint Review qua Google Meet, trình diễn sản phẩm và tổng kết quá trình làm việc. Hiển thực hiện phần trình diễn; Hiệp đóng vai người nhận bàn giao.


Nhóm đã tái dựng biểu đồ Burndown theo số task còn lại trong từng Sprint từ lịch sử đóng và mở lại Issue trên GitHub. Biểu đồ phục vụ theo dõi tiến độ; đơn vị là số task, không phải điểm ước lượng công việc hoặc số commit.
\vspace{1em}
\textbf{Biểu đồ Burndown theo Story Points}

Điểm công việc được ước lượng hồi cứu theo độ phức tạp tương đối của từng task (mỗi task 1 điểm theo đề). Tiến độ được tái dựng từ ngày bàn giao từng task trong Chương 2, chốt cuối ngày 02/11/2026. Đường lý tưởng giảm đều đến hết Sprint; đường thực tế thể hiện số điểm của các task chưa hoàn tất. Ba biểu đồ dưới đây được vẽ trực tiếp bằng TikZ/pgfplots.

\begin{center}
\begin{tikzpicture}
\begin{axis}[
  width=0.92\linewidth, height=6cm,
  title={Sprint 1 (14–20/10/2026): T01–T04},
  xlabel={Ngày}, ylabel={Điểm còn lại},
  ymin=0, ymax=4.4, xmin=-0.3, xmax=6.3,
  xtick={0,1,2,3,4,5,6},
  xticklabels={14/10,15/10,16/10,17/10,18/10,19/10,20/10},
  x tick label style={font=\scriptsize},
  ytick={0,1,2,3,4},
  legend style={at={(0.98,0.98)}, anchor=north east, font=\scriptsize},
  grid=major, grid style={gray!20},
  title style={font=\small\bfseries},
]
\addplot[thick, gray, dashed, mark=none] coordinates {(0,4)(1,3.33)(2,2.67)(3,2)(4,1.33)(5,0.67)(6,0)};
\addlegendentry{Lý tưởng}
\addplot[thick, blue!70!black, mark=*, mark size=1.5pt] coordinates {(0,3)(1,2)(2,2)(3,1)(4,0)(5,0)(6,0)};
\addlegendentry{Thực tế}
\end{axis}
\end{tikzpicture}

\vspace{0.8em}
\begin{tikzpicture}
\begin{axis}[
  width=0.92\linewidth, height=6cm,
  title={Sprint 2 (21–27/10/2026): T05–T07},
  xlabel={Ngày}, ylabel={Điểm còn lại},
  ymin=0, ymax=3.4, xmin=-0.3, xmax=6.3,
  xtick={0,1,2,3,4,5,6},
  xticklabels={21/10,22/10,23/10,24/10,25/10,26/10,27/10},
  x tick label style={font=\scriptsize},
  ytick={0,1,2,3},
  legend style={at={(0.98,0.98)}, anchor=north east, font=\scriptsize},
  grid=major, grid style={gray!20},
  title style={font=\small\bfseries},
]
\addplot[thick, gray, dashed, mark=none] coordinates {(0,3)(1,2.5)(2,2)(3,1.5)(4,1)(5,0.5)(6,0)};
\addlegendentry{Lý tưởng}
\addplot[thick, blue!70!black, mark=*, mark size=1.5pt] coordinates {(0,2)(1,2)(2,1)(3,1)(4,0)(5,0)(6,0)};
\addlegendentry{Thực tế}
\end{axis}
\end{tikzpicture}

\vspace{0.8em}
\begin{tikzpicture}
\begin{axis}[
  width=0.92\linewidth, height=6cm,
  title={Sprint 3 (28/10–03/11/2026): T08–T10},
  xlabel={Ngày}, ylabel={Điểm còn lại},
  ymin=0, ymax=3.4, xmin=-0.3, xmax=6.3,
  xtick={0,1,2,3,4,5,6},
  xticklabels={28/10,29/10,30/10,31/10,01/11,02/11,03/11},
  x tick label style={font=\scriptsize},
  ytick={0,1,2,3},
  legend style={at={(0.98,0.98)}, anchor=north east, font=\scriptsize},
  grid=major, grid style={gray!20},
  title style={font=\small\bfseries},
]
\addplot[thick, gray, dashed, mark=none] coordinates {(0,3)(1,2.5)(2,2)(3,1.5)(4,1)(5,0.5)(6,0)};
\addlegendentry{Lý tưởng}
\addplot[thick, blue!70!black, mark=*, mark size=1.5pt] coordinates {(0,3)(1,2)(2,2)(3,2)(4,2)(5,1)(6,0)};
\addlegendentry{Thực tế}
\end{axis}
\end{tikzpicture}
\end{center}


\chapter{NHẬT KÝ TIẾN ĐỘ VÀ CHI TIẾT THỰC HIỆN CÁC TASK}

Nhật ký dưới đây ghi nhận phân công, phần việc bổ sung, chỉnh sửa, kiểm tra và bàn giao của từng task. Nội dung các bảng được giữ theo bản nhóm đã tổng hợp; thời gian tính theo giờ Việt Nam. Người thực hiện bàn giao phần việc cho nhóm, còn việc tích hợp và nghiệm thu là các bước khác nhau.


\Needspace{8\baselineskip}
\section{Task 1: Khởi tạo dự án và Kiến trúc nền tảng (T01)}

\begingroup
\small
\setlength{\tabcolsep}{5pt}
\renewcommand{\arraystretch}{1.2}
\begin{longtable}{|>{\centering\arraybackslash}p{\dimexpr 0.07785\linewidth-2\tabcolsep-2\arrayrulewidth\relax}|>{\raggedright\arraybackslash}p{\dimexpr 0.18559\linewidth-2\tabcolsep-2\arrayrulewidth\relax}|>{\raggedright\arraybackslash}p{\dimexpr 0.14968\linewidth-2\tabcolsep-2\arrayrulewidth\relax}|>{\raggedright\arraybackslash}p{\dimexpr 0.58688\linewidth-2\tabcolsep-2\arrayrulewidth\relax}|}
\hline
\textbf{STT} & \textbf{Thời gian} & \multicolumn{2}{|>{\raggedright\arraybackslash}p{\dimexpr 0.73656\linewidth-2\tabcolsep-2\arrayrulewidth\relax}|}{\textbf{Công việc}} \\ \hline
\endfirsthead
\hline
\textbf{STT} & \textbf{Thời gian} & \multicolumn{2}{|>{\raggedright\arraybackslash}p{\dimexpr 0.73656\linewidth-2\tabcolsep-2\arrayrulewidth\relax}|}{\textbf{Công việc}} \\ \hline
\endhead
\endfoot
\endlastfoot
1 &  & Phân công & Hiển: khởi tạo dự án Next.js/NestJS, Git và môi trường chạy. \\ \hline
2 & 14/10/2026 & Bổ sung & Hiển khởi tạo kho mã nguồn và hai ứng dụng giao diện, máy chủ; thiết lập tệp môi trường mẫu và hướng dẫn chạy. \\ \hline
3 & 14/10/2026 & Bàn giao & Hiển bàn giao cấu trúc dự án Next.js/NestJS, cấu hình chạy và hướng dẫn khởi động cho các thành viên phát triển tiếp. \\ \hline
\end{longtable}
\endgroup


\Needspace{8\baselineskip}
\section{Task 2: Thiết kế Cơ sở dữ liệu và API Sản phẩm (T02)}

\begingroup
\small
\setlength{\tabcolsep}{5pt}
\renewcommand{\arraystretch}{1.2}
\begin{longtable}{|>{\centering\arraybackslash}p{\dimexpr 0.07785\linewidth-2\tabcolsep-2\arrayrulewidth\relax}|>{\raggedright\arraybackslash}p{\dimexpr 0.18559\linewidth-2\tabcolsep-2\arrayrulewidth\relax}|>{\raggedright\arraybackslash}p{\dimexpr 0.14968\linewidth-2\tabcolsep-2\arrayrulewidth\relax}|>{\raggedright\arraybackslash}p{\dimexpr 0.58688\linewidth-2\tabcolsep-2\arrayrulewidth\relax}|}
\hline
\textbf{STT} & \textbf{Thời gian} & \multicolumn{2}{|>{\raggedright\arraybackslash}p{\dimexpr 0.73656\linewidth-2\tabcolsep-2\arrayrulewidth\relax}|}{\textbf{Công việc}} \\ \hline
\endfirsthead
\hline
\textbf{STT} & \textbf{Thời gian} & \multicolumn{2}{|>{\raggedright\arraybackslash}p{\dimexpr 0.73656\linewidth-2\tabcolsep-2\arrayrulewidth\relax}|}{\textbf{Công việc}} \\ \hline
\endhead
\endfoot
\endlastfoot
1 & 14/10/2026 & Phân công & Thành: cơ sở dữ liệu Prisma/PostgreSQL và API sản phẩm. \\ \hline
2 & 15/10/2026 & Bổ sung & Thành tạo cấu trúc dữ liệu sản phẩm, dữ liệu mẫu và hai chức năng lấy danh sách, xem chi tiết đồ uống từ cơ sở dữ liệu. \\ \hline
3 & 15/10/2026 & Bàn giao & Thành bàn giao API sản phẩm, cơ sở dữ liệu và hướng dẫn sử dụng để nhóm kết nối với giao diện Menu. \\ \hline
4 & 03/11/2026 & Chỉnh sửa & Hiển bổ sung kiểm tra mã sản phẩm khi xem chi tiết: từ chối mã sai định dạng hoặc vượt giới hạn trước khi truy vấn dữ liệu, trả thông báo dữ liệu không hợp lệ thay vì lỗi máy chủ; mã hợp lệ nhưng không có sản phẩm vẫn báo không tìm thấy. \\ \hline
\end{longtable}
\endgroup


\Needspace{8\baselineskip}
\section{Task 3: Phát triển Giao diện Thực đơn (T03)}

\begingroup
\small
\setlength{\tabcolsep}{5pt}
\renewcommand{\arraystretch}{1.2}
\begin{longtable}{|>{\centering\arraybackslash}p{\dimexpr 0.07785\linewidth-2\tabcolsep-2\arrayrulewidth\relax}|>{\raggedright\arraybackslash}p{\dimexpr 0.18559\linewidth-2\tabcolsep-2\arrayrulewidth\relax}|>{\raggedright\arraybackslash}p{\dimexpr 0.14968\linewidth-2\tabcolsep-2\arrayrulewidth\relax}|>{\raggedright\arraybackslash}p{\dimexpr 0.58688\linewidth-2\tabcolsep-2\arrayrulewidth\relax}|}
\hline
\textbf{STT} & \textbf{Thời gian} & \multicolumn{2}{|>{\raggedright\arraybackslash}p{\dimexpr 0.73656\linewidth-2\tabcolsep-2\arrayrulewidth\relax}|}{\textbf{Công việc}} \\ \hline
\endfirsthead
\hline
\textbf{STT} & \textbf{Thời gian} & \multicolumn{2}{|>{\raggedright\arraybackslash}p{\dimexpr 0.73656\linewidth-2\tabcolsep-2\arrayrulewidth\relax}|}{\textbf{Công việc}} \\ \hline
\endhead
\endfoot
\endlastfoot
1 & 14/10/2026 & Phân công & Xuyến: Menu, dữ liệu API và ảnh sản phẩm. \\ \hline
2 & 16/10/2026 & Bổ sung & Xuyến xây dựng Menu hiển thị tên, ảnh và giá đồ uống từ API; xử lý lúc tải dữ liệu, danh sách rỗng và lỗi kết nối. \\ \hline
3 & 17/10/2026 & Chỉnh sửa & Xuyến thay ảnh sản phẩm, cập nhật hướng dẫn và bàn giao giao diện Menu sau khi chỉnh sửa. \\ \hline
4 & 17/10/2026 & Bàn giao & Xuyến bàn giao giao diện Menu đã bổ sung ảnh đồ uống, kết nối dữ liệu sản phẩm và xử lý trạng thái tải, rỗng, lỗi. \\ \hline
\end{longtable}
\endgroup


\Needspace{8\baselineskip}
\section{Task 4: Màn hình Chi tiết Sản phẩm \& Tùy chọn (T04)}

\begingroup
\small
\setlength{\tabcolsep}{5pt}
\renewcommand{\arraystretch}{1.2}
\begin{longtable}{|>{\centering\arraybackslash}p{\dimexpr 0.07785\linewidth-2\tabcolsep-2\arrayrulewidth\relax}|>{\raggedright\arraybackslash}p{\dimexpr 0.18559\linewidth-2\tabcolsep-2\arrayrulewidth\relax}|>{\raggedright\arraybackslash}p{\dimexpr 0.14968\linewidth-2\tabcolsep-2\arrayrulewidth\relax}|>{\raggedright\arraybackslash}p{\dimexpr 0.58688\linewidth-2\tabcolsep-2\arrayrulewidth\relax}|}
\hline
\textbf{STT} & \textbf{Thời gian} & \multicolumn{2}{|>{\raggedright\arraybackslash}p{\dimexpr 0.73656\linewidth-2\tabcolsep-2\arrayrulewidth\relax}|}{\textbf{Công việc}} \\ \hline
\endfirsthead
\hline
\textbf{STT} & \textbf{Thời gian} & \multicolumn{2}{|>{\raggedright\arraybackslash}p{\dimexpr 0.73656\linewidth-2\tabcolsep-2\arrayrulewidth\relax}|}{\textbf{Công việc}} \\ \hline
\endhead
\endfoot
\endlastfoot
1 & 14/10/2026 & Phân công & Hiệp: trang chi tiết, kích cỡ và topping. \\ \hline
2 & 18/10/2026 & Bổ sung & Hiệp xây dựng trang chi tiết đồ uống; cho chọn kích cỡ S/M/L, topping và cập nhật giá theo lựa chọn. \\ \hline
3 & 18/10/2026 & Bàn giao & Hiệp bàn giao trang chi tiết cùng chức năng chọn kích cỡ, topping và tính giá để kết nối với giỏ hàng. \\ \hline
\end{longtable}
\endgroup


\Needspace{8\baselineskip}
\section{Task 5: Quản lý Giỏ hàng linh hoạt (T05)}

\begingroup
\small
\setlength{\tabcolsep}{5pt}
\renewcommand{\arraystretch}{1.2}
\begin{longtable}{|>{\centering\arraybackslash}p{\dimexpr 0.07785\linewidth-2\tabcolsep-2\arrayrulewidth\relax}|>{\raggedright\arraybackslash}p{\dimexpr 0.18559\linewidth-2\tabcolsep-2\arrayrulewidth\relax}|>{\raggedright\arraybackslash}p{\dimexpr 0.14968\linewidth-2\tabcolsep-2\arrayrulewidth\relax}|>{\raggedright\arraybackslash}p{\dimexpr 0.58688\linewidth-2\tabcolsep-2\arrayrulewidth\relax}|}
\hline
\textbf{STT} & \textbf{Thời gian} & \multicolumn{2}{|>{\raggedright\arraybackslash}p{\dimexpr 0.73656\linewidth-2\tabcolsep-2\arrayrulewidth\relax}|}{\textbf{Công việc}} \\ \hline
\endfirsthead
\hline
\textbf{STT} & \textbf{Thời gian} & \multicolumn{2}{|>{\raggedright\arraybackslash}p{\dimexpr 0.73656\linewidth-2\tabcolsep-2\arrayrulewidth\relax}|}{\textbf{Công việc}} \\ \hline
\endhead
\endfoot
\endlastfoot
1 & 14/10/2026 & Phân công & Thiện: giỏ hàng, số lượng và tổng tiền. \\ \hline
2 & 20/10/2026 & Bổ sung & Thiện xây dựng giỏ hàng: thêm món, tăng giảm số lượng, xóa món, hiển thị số lượng trên biểu tượng giỏ và tính tổng tiền. \\ \hline
3 & 21/10/2026 & Chỉnh sửa & Thiện sửa bố cục, nhãn tiếng Việt và điều hướng về Menu sau khi thêm món; bổ sung hướng dẫn sử dụng giỏ hàng. \\ \hline
4 & 21/10/2026 & Bàn giao & Thiện bàn giao giỏ hàng đã chỉnh sửa, gồm quản lý món, số lượng và tổng tiền để phát triển bước thanh toán. \\ \hline
\end{longtable}
\endgroup


\Needspace{8\baselineskip}
\section{Task 6: API Tạo đơn hàng và Tính toán Giá (T06)}

\begingroup
\small
\setlength{\tabcolsep}{5pt}
\renewcommand{\arraystretch}{1.2}
\begin{longtable}{|>{\centering\arraybackslash}p{\dimexpr 0.07785\linewidth-2\tabcolsep-2\arrayrulewidth\relax}|>{\raggedright\arraybackslash}p{\dimexpr 0.18559\linewidth-2\tabcolsep-2\arrayrulewidth\relax}|>{\raggedright\arraybackslash}p{\dimexpr 0.14968\linewidth-2\tabcolsep-2\arrayrulewidth\relax}|>{\raggedright\arraybackslash}p{\dimexpr 0.58688\linewidth-2\tabcolsep-2\arrayrulewidth\relax}|}
\hline
\textbf{STT} & \textbf{Thời gian} & \multicolumn{2}{|>{\raggedright\arraybackslash}p{\dimexpr 0.73656\linewidth-2\tabcolsep-2\arrayrulewidth\relax}|}{\textbf{Công việc}} \\ \hline
\endfirsthead
\hline
\textbf{STT} & \textbf{Thời gian} & \multicolumn{2}{|>{\raggedright\arraybackslash}p{\dimexpr 0.73656\linewidth-2\tabcolsep-2\arrayrulewidth\relax}|}{\textbf{Công việc}} \\ \hline
\endhead
\endfoot
\endlastfoot
1 & 14/10/2026 & Phân công & Hiển: API tạo đơn, kiểm tra dữ liệu và tính giá. \\ \hline
2 & 23/10/2026 & Bổ sung & Hiển xây dựng chức năng tạo đơn: kiểm tra món, kích cỡ, topping và số lượng; tính lại giá tại máy chủ rồi lưu đơn ở trạng thái chờ thanh toán. \\ \hline
3 & 23/10/2026 & Bàn giao & Hiển bàn giao API tạo đơn và hướng dẫn kiểm thử để nhóm kết nối với đăng nhập, giỏ hàng và thanh toán. \\ \hline
4 & 03/11/2026 & Chỉnh sửa & Hiển kiểm tra giới hạn mã sản phẩm trong yêu cầu tạo đơn và mã đơn khi tra cứu; chặn dữ liệu không hợp lệ trước khi truy vấn cơ sở dữ liệu, giữ cách xử lý không tìm thấy đối với mã hợp lệ. \\ \hline
\end{longtable}
\endgroup


\Needspace{8\baselineskip}
\section{Task 7: Mô-đun Xác thực Người dùng \& Bảo mật (T07)}

\begingroup
\small
\setlength{\tabcolsep}{5pt}
\renewcommand{\arraystretch}{1.2}
\begin{longtable}{|>{\centering\arraybackslash}p{\dimexpr 0.07785\linewidth-2\tabcolsep-2\arrayrulewidth\relax}|>{\raggedright\arraybackslash}p{\dimexpr 0.18559\linewidth-2\tabcolsep-2\arrayrulewidth\relax}|>{\raggedright\arraybackslash}p{\dimexpr 0.14968\linewidth-2\tabcolsep-2\arrayrulewidth\relax}|>{\raggedright\arraybackslash}p{\dimexpr 0.58688\linewidth-2\tabcolsep-2\arrayrulewidth\relax}|}
\hline
\textbf{STT} & \textbf{Thời gian} & \multicolumn{2}{|>{\raggedright\arraybackslash}p{\dimexpr 0.73656\linewidth-2\tabcolsep-2\arrayrulewidth\relax}|}{\textbf{Công việc}} \\ \hline
\endfirsthead
\hline
\textbf{STT} & \textbf{Thời gian} & \multicolumn{2}{|>{\raggedright\arraybackslash}p{\dimexpr 0.73656\linewidth-2\tabcolsep-2\arrayrulewidth\relax}|}{\textbf{Công việc}} \\ \hline
\endhead
\endfoot
\endlastfoot
1 & 14/10/2026 & Phân công & Khoa: đăng ký/đăng nhập, JWT và thông báo xác thực. \\ \hline
2 & 25/10/2026 & Bàn giao & Khoa bàn giao đăng ký, đăng nhập và bảo vệ API bằng mã xác thực; mật khẩu được băm, đơn hàng gắn với tài khoản và có kiểm thử đăng xuất/chuyển tài khoản. \\ \hline
3 & 30/10/2026 & Kiểm tra & Hiển kiểm tra lại đăng ký, đăng nhập và quyền truy cập API; bổ sung kết quả kiểm thử vào hướng dẫn dự án. \\ \hline
4 & 02/11/2026 & Chỉnh sửa & Khoa sửa lỗi hiển thị ngay dưới ô email hoặc mật khẩu; phân biệt sai thông tin đăng nhập với mất kết nối và bàn giao bản sửa kèm kiểm thử. \\ \hline
5 & 02/11/2026 & Bàn giao & Khoa bàn giao bản sửa đăng ký, đăng nhập và thông báo lỗi theo ô nhập, kèm kết quả kiểm thử để nhóm sử dụng bản hoàn thiện. \\ \hline
\end{longtable}
\endgroup


\Needspace{8\baselineskip}
\section{Task 8: Mô phỏng Thanh toán Trực tuyến (T08)}

\begingroup
\small
\setlength{\tabcolsep}{5pt}
\renewcommand{\arraystretch}{1.2}
\begin{longtable}{|>{\centering\arraybackslash}p{\dimexpr 0.07785\linewidth-2\tabcolsep-2\arrayrulewidth\relax}|>{\raggedright\arraybackslash}p{\dimexpr 0.18559\linewidth-2\tabcolsep-2\arrayrulewidth\relax}|>{\raggedright\arraybackslash}p{\dimexpr 0.14968\linewidth-2\tabcolsep-2\arrayrulewidth\relax}|>{\raggedright\arraybackslash}p{\dimexpr 0.58688\linewidth-2\tabcolsep-2\arrayrulewidth\relax}|}
\hline
\textbf{STT} & \textbf{Thời gian} & \textbf{Công việc} & \textbf{Nội dung} \\ \hline
\endfirsthead
\hline
\textbf{STT} & \textbf{Thời gian} & \textbf{Công việc} & \textbf{Nội dung} \\ \hline
\endhead
\endfoot
\endlastfoot
1 & 14/10/2026 & Phân công & Thành phụ trách chức năng thanh toán. \\ \hline
2 & 25/10/2026 & Bổ sung phân công & Thiện phụ trách giao diện checkout và thử lại thanh toán; Thành phụ trách API thanh toán. \\ \hline
3 & 28/10/2026 & Bàn giao & Thành bàn giao API thanh toán ví/thẻ mô phỏng; xử lý thành công, thất bại và thử lại, kèm kiểm thử kết quả. \\ \hline
4 & 28–29/10/2026 & Bàn giao & Thiện bàn giao giao diện xác nhận thông tin đơn, chọn cách thanh toán và thử lại khi lỗi. Thành tổng hợp giao diện với API thành luồng thanh toán hoàn chỉnh. \\ \hline
5 & 29/10/2026 & Chỉnh sửa & Thiện sửa việc lưu đơn chờ theo từng tài khoản và xử lý giao diện sau góp ý; bàn giao bản checkout đã chỉnh sửa. \\ \hline
6 & 01/11/2026 & Chỉnh sửa & Hiển bổ sung số tiền phải trả ngay trên nút thanh toán để khách kiểm tra trước khi xác nhận. \\ \hline
7 & 03/11/2026 & Chỉnh sửa & Hiển bổ sung giới hạn mã đơn gửi lên API thanh toán; yêu cầu có mã vượt phạm vi được từ chối trước khi xử lý, tránh lỗi máy chủ. Giữ nguyên luồng thanh toán thành công, thất bại và thử lại. \\ \hline
\end{longtable}
\endgroup


\Needspace{8\baselineskip}
\section{Task 9: Màn hình Xác nhận, Lịch sử \& Đóng gói Docker (T09)}

\begingroup
\small
\setlength{\tabcolsep}{5pt}
\renewcommand{\arraystretch}{1.2}
\begin{longtable}{|>{\centering\arraybackslash}p{\dimexpr 0.07785\linewidth-2\tabcolsep-2\arrayrulewidth\relax}|>{\raggedright\arraybackslash}p{\dimexpr 0.18559\linewidth-2\tabcolsep-2\arrayrulewidth\relax}|>{\raggedright\arraybackslash}p{\dimexpr 0.14968\linewidth-2\tabcolsep-2\arrayrulewidth\relax}|>{\raggedright\arraybackslash}p{\dimexpr 0.58688\linewidth-2\tabcolsep-2\arrayrulewidth\relax}|}
\hline
\textbf{STT} & \textbf{Thời gian} & \multicolumn{2}{|>{\raggedright\arraybackslash}p{\dimexpr 0.73656\linewidth-2\tabcolsep-2\arrayrulewidth\relax}|}{\textbf{Công việc}} \\ \hline
\endfirsthead
\hline
\textbf{STT} & \textbf{Thời gian} & \multicolumn{2}{|>{\raggedright\arraybackslash}p{\dimexpr 0.73656\linewidth-2\tabcolsep-2\arrayrulewidth\relax}|}{\textbf{Công việc}} \\ \hline
\endhead
\endfoot
\endlastfoot
1 & 14/10/2026 & Phân công & Hiệp và Xuyến \\ \hline
2 & 28/10/2026 & Bổ sung phân công & Xuyến: xác nhận và lịch sử; Hiệp: Docker và hướng dẫn chạy. \\ \hline
3 & 31/10/2026 & Bàn giao & Xuyến xây dựng và bàn giao trang xác nhận, danh sách lịch sử và chi tiết đơn; hiển thị mã đơn, trạng thái cùng số tiền. \\ \hline
4 & 31/10/2026 & Bổ sung & Hiệp bổ sung và bàn giao bộ chạy Docker gồm giao diện, API, PostgreSQL cùng hướng dẫn cài đặt và nạp dữ liệu mẫu. \\ \hline
5 & 01/11/2026 & Bàn giao & Hiệp sửa việc giữ thông tin xác nhận sau thanh toán và lỗi định dạng giao diện khi đóng gói; bàn giao bản Docker và hướng dẫn đã chỉnh sửa. \\ \hline
6 & 03/11/2026 & Chỉnh sửa & Hiển bổ sung tùy chọn mạng khi build Docker và cho bước cài thư viện backend thử lại tối đa ba lần khi tải bị ngắt; nếu vẫn lỗi thì dừng build. Cập nhật README với lệnh chạy, tài khoản ví dụ và hướng dẫn kiểm thử. \\ \hline
7 & 03/11/2026 & Kiểm tra và bàn giao & Hiển kiểm tra bản Docker trên môi trường riêng: PostgreSQL và API hoạt động ổn định, giao diện khởi động được; 35 kiểm tra HTTP và 21 kiểm tra giao diện đạt. Bàn giao bản sửa cùng hướng dẫn đã cập nhật; tích hợp vào main ngày 03/11/2026. \\ \hline
\end{longtable}
\endgroup


\Needspace{8\baselineskip}
\section{Task 10: Xử lý Nghiệp vụ Backend Nâng cao (T10)}

\begingroup
\small
\setlength{\tabcolsep}{5pt}
\renewcommand{\arraystretch}{1.2}
\begin{longtable}{|>{\centering\arraybackslash}p{\dimexpr 0.07785\linewidth-2\tabcolsep-2\arrayrulewidth\relax}|>{\raggedright\arraybackslash}p{\dimexpr 0.18559\linewidth-2\tabcolsep-2\arrayrulewidth\relax}|>{\raggedright\arraybackslash}p{\dimexpr 0.14968\linewidth-2\tabcolsep-2\arrayrulewidth\relax}|>{\raggedright\arraybackslash}p{\dimexpr 0.58688\linewidth-2\tabcolsep-2\arrayrulewidth\relax}|}
\hline
\textbf{STT} & \textbf{Thời gian} & \multicolumn{2}{|>{\raggedright\arraybackslash}p{\dimexpr 0.73656\linewidth-2\tabcolsep-2\arrayrulewidth\relax}|}{\textbf{Công việc}} \\ \hline
\endfirsthead
\hline
\textbf{STT} & \textbf{Thời gian} & \multicolumn{2}{|>{\raggedright\arraybackslash}p{\dimexpr 0.73656\linewidth-2\tabcolsep-2\arrayrulewidth\relax}|}{\textbf{Công việc}} \\ \hline
\endhead
\endfoot
\endlastfoot
1 & 14/10/2026 & Phân công & Khoa và Hiển \\ \hline
2 & 28/10/2026 & Bổ sung phân công & Khoa: T10A trạng thái, khuyến mãi, điểm; Hiển: T10B chống lặp, tồn kho và nhân viên. \\ \hline
3 & 30/10/2026 & Bổ sung & Khoa xây dựng trạng thái đơn, mã giảm giá và điểm thưởng; bổ sung dữ liệu, quy tắc xử lý và tài liệu giải thích. \\ \hline
4 & 31/10/2026 & Bàn giao & Khoa bàn giao trạng thái đơn, khuyến mãi, điểm và kiểm thử. Hiển bàn giao cơ chế chống thanh toán lặp và trừ kho trong giao dịch, kiểm soát nhiều đơn cùng mua để tránh âm kho. \\ \hline
5 & 01/11/2026 & Bổ sung & Hiển bổ sung trang quản lý quầy: xem nguyên liệu, điều chỉnh kho có lý do, theo dõi đơn và điểm; phân quyền nhân viên với khách hàng. \\ \hline
6 & 02/11/2026 & Chỉnh sửa & Hiển sửa kiểm tra đủ nguyên liệu theo kích cỡ, số phần topping trước khi thêm món; hủy đơn đã thanh toán trước khi pha sẽ hoàn kho và thu hồi điểm đúng một lần. Bàn giao bản tổng hợp cùng ảnh, log kiểm thử. \\ \hline
7 & 02/11/2026 & Bàn giao & Khoa bàn giao phần trạng thái đơn, khuyến mãi và điểm; Hiển tổng hợp với chống thanh toán lặp, tồn kho, trang nhân viên và bản sửa hủy đơn; bàn giao mã nguồn cùng ảnh, log kiểm thử cho nhóm. \\ \hline
8 & 03/11/2026 & Chỉnh sửa & Hiển áp dụng kiểm tra mã nguyên liệu và mã đơn trên API nhân viên; bổ sung giới hạn mã sản phẩm cho bước kiểm tra tồn kho của giỏ. Chặn mã không hợp lệ trước khi truy vấn, giữ nguyên quy tắc trừ kho, hoàn kho và chống thanh toán lặp. \\ \hline
9 & 03/11/2026 & Kiểm tra & Hiển bổ sung 58 trường hợp kiểm thử hồi quy cho mã không hợp lệ, giá trị biên và mã không tồn tại trên các API liên quan. Bốn bộ kiểm thử ghi nhận 248/248 đạt; kiểm tra mã nguồn và build hai ứng dụng đều đạt trước khi bàn giao. \\ \hline
\end{longtable}
\endgroup


Nhật ký các task được đối chiếu với lịch sử mã nguồn. Kết quả Sprint Review ngày 03/11/2026 và minh chứng buổi trình diễn được tổng hợp tại Issue \#11.


\Needspace{8\baselineskip}
\section{Tổng hợp kết quả phát triển}

Nhóm đã tích hợp các phần nền tảng, sản phẩm, Menu, chi tiết, giỏ hàng, tạo đơn, xác thực, thanh toán, xác nhận và lịch sử thành luồng mua hàng. Các bảng T01–T10 phía trên giữ vai trò nhật ký công việc; kết quả kiểm thử được trình bày riêng ở chương tiếp theo.


Phần nghiệp vụ nâng cao gồm kiểm soát chuyển trạng thái, xử lý yêu cầu thanh toán lặp, kiểm soát kho khi nhiều yêu cầu cùng thực hiện, mã giảm giá và điểm thưởng. Đơn đã thanh toán chỉ được hủy trước khi bắt đầu pha; bản sửa xử lý hoàn kho và thu hồi điểm đúng một lần. Các bước chuyển không hợp lệ bị từ chối.


Qua quá trình tích hợp, nhóm đã sửa lỗi định dạng giao diện ở bản đóng gói, lưu thông tin xác nhận sau thanh toán và thông báo xác thực. Trang nhân viên được tách khỏi chức năng mua hàng; nhân viên xem điểm và lịch sử của khách, không điều chỉnh điểm tùy ý.


\chapter{KIỂM THỬ PHẦN MỀM VÀ HƯỚNG DẪN BÀN GIAO}

Nhóm sử dụng kiểm thử đơn vị để kiểm tra hàm và trạng thái, kiểm thử API để kiểm tra dữ liệu và quyền truy cập, kiểm thử PostgreSQL để kiểm tra giao dịch và đồng thời, cùng thao tác trình duyệt để đối chiếu trải nghiệm thực tế. Lint và build được dùng để kiểm tra tính hợp lệ của mã nguồn và khả năng tạo bản chạy. Bảng sau tổng hợp kết quả có ngày, phiên bản và phạm vi ghi nhận; các lượt kiểm tra độc lập không được cộng trùng.


\begingroup
\small
\setlength{\tabcolsep}{5pt}
\renewcommand{\arraystretch}{1.2}
\begin{longtable}{|>{\raggedright\arraybackslash}p{\dimexpr 0.27545\linewidth-2\tabcolsep-2\arrayrulewidth\relax}|>{\raggedright\arraybackslash}p{\dimexpr 0.72455\linewidth-2\tabcolsep-2\arrayrulewidth\relax}|}
\hline
\textbf{Mốc kiểm tra} & \textbf{Kết quả ghi nhận} \\ \hline
\endfirsthead
\hline
\textbf{Mốc kiểm tra} & \textbf{Kết quả ghi nhận} \\ \hline
\endhead
\endfoot
\endlastfoot
02/11 • fe832a5 / 005d1ef & Backend unit 82/82; E2E 44/44; PostgreSQL integration 27/27; frontend unit 30/30. Tổng 183/183; lint/build hai ứng dụng đạt. \\ \hline
02/11 • 0730b95 & Frontend unit 37/37, lint/build đạt; kiểm tra lại thông báo lỗi xác thực trên trình duyệt với API giả lập. Đã merge thành main ddb3013. \\ \hline
02/11 • giao diện và DB thật & Playwright kết nối API/PostgreSQL thử riêng: chặn vượt kho; thanh toán 2 ly 72.000đ, cộng 7 điểm; hủy PAID hoàn kho và thu hồi điểm. Có 6 ảnh và log. \\ \hline
02/11 • đo API trên 005d1ef & 10 API \ensuremath{\times} 30 mẫu: 300/300 dưới 500 ms; lớn nhất 102,14 ms. Đo cục bộ sau khởi động nóng, tuần tự; không đại diện tải cao hoặc mạng Internet. \\ \hline
03/11 • bản sửa bàn giao & Backend unit 82/82; frontend unit 37/37; E2E 102/102; PostgreSQL integration 27/27. Tổng 248/248. Bản Docker ghi nhận 35/35 kiểm tra HTTP và 21/21 kiểm tra giao diện đạt. \\ \hline
\end{longtable}
\endgroup


\Needspace{8\baselineskip}
\section{Phân tích kết quả kiểm thử}

Các con số trong bảng thuộc từng bản mã nguồn và từng lần chạy đã ghi nhận. Lượt kiểm tra ngày 03/11 sau khi sửa giới hạn mã dữ liệu ghi nhận 248/248 bài đạt; các con số ngày 02/11 được giữ làm lịch sử, không cộng dồn với lượt chạy mới. Lượt kiểm tra tích hợp ngày 02/11 có tổng 183 bài đạt, gồm kiểm thử đơn vị, API và PostgreSQL. Sau đó, phần sửa giao diện xác thực được kiểm tra riêng với 37 bài frontend đạt; không cộng hai lần chạy thành một tổng mới.


Minh chứng giao diện T10B sử dụng trình duyệt do Playwright điều khiển, kết nối API và PostgreSQL thử riêng. Với 70 g thạch đủ hai phần, giỏ chặn tăng quá hai ly; đơn hai ly 72.000đ thanh toán thành công được cộng 7 điểm. Kho giảm theo công thức; khi hủy trước lúc pha, nguyên liệu được hoàn về mức ban đầu và điểm trở về 0.


Bộ kiểm thử dữ liệu kiểm tra yêu cầu trùng, nhiều yêu cầu tranh kho, lỗi giữa giao dịch và hủy đơn đồng thời. Kết quả này chứng minh các trường hợp đã xây dựng; không có nghĩa mọi tình huống vận hành đều đã được kiểm tra.


Kết quả đo 10 API, mỗi API 30 mẫu sau các lần làm nóng, ghi nhận cả 300 mẫu dưới 500 ms; mẫu lớn nhất là 102,14 ms. Phép đo chạy tuần tự trên máy cục bộ và dữ liệu thử, chưa đại diện cho tải cao hoặc đường truyền Internet. Riêng kiểm tra phân loại lỗi giao diện xác thực có sử dụng phản hồi API giả lập để tái hiện lỗi.


\Needspace{8\baselineskip}
\section{Hướng dẫn chạy và kiểm tra hệ thống}

Máy chạy cần có Git và Docker Desktop. Lấy mã nguồn từ kho Git của nhóm, mở terminal tại thư mục gốc chứa docker-compose.yml, sao chép .env.example thành .env và điền cấu hình theo README, gồm khóa JWT riêng cho môi trường chạy. Khởi động Docker Desktop, đợi Engine running rồi thực hiện:


\begin{quote}\ttfamily\small docker compose up --build -d\end{quote}

\begin{quote}\ttfamily\small docker compose ps\end{quote}

Kiểm tra PostgreSQL và API có trạng thái healthy, giao diện đã chạy. Lần đầu nạp dữ liệu mẫu bằng lệnh docker compose exec api npm run prisma:seed. Khách mở http://localhost:3000; nhân viên mở http://localhost:3000/staff bằng tài khoản được cấp quyền. Dữ liệu được lưu trong volume PostgreSQL.


Những lần sau dùng docker compose up -d; sau khi thay đổi mã nguồn cần chạy lại lệnh build ở trên. Dừng bằng docker compose down; tránh tùy chọn -v nếu cần giữ dữ liệu. Khi có lỗi, kiểm tra trạng thái dịch vụ và log API, đồng thời đối chiếu URL API, cổng và cấu hình môi trường trong README.


Để chạy kiểm thử tự động, cài thư viện trong từng ứng dụng rồi thực hiện npm run lint, npm test và npm run build. Với API/E2E và bài PostgreSQL, dùng cơ sở dữ liệu thử riêng, cấu hình TEST\_DATABASE\_URL, chạy generate, migration và seed trước khi chạy npm run test:e2e và npm run test:database tại backend. Không dùng dữ liệu đang trình diễn cho các bài test có ghi hoặc xóa dữ liệu.


Khi chạy thử giao diện, kiểm tra giá theo tùy chọn, tổng giỏ, quyền đăng nhập, thành công/thất bại thanh toán và lịch sử đơn. Nếu có lỗi, ghi ngày, bản mã nguồn, bước tái hiện, kết quả mong đợi, kết quả thực tế và ảnh/log; sau khi sửa cần kiểm tra lại trường hợp đó.


\Needspace{8\baselineskip}
\section{Tổng kết quá trình thực hiện}

Qua quá trình phát triển, nhóm đã có bản chạy tích hợp và kết quả kiểm thử cho các chức năng A1. Những khó khăn đã ghi nhận gồm giao diện production, lưu xác nhận sau thanh toán, thống nhất quy tắc hủy và hoàn kho, cùng cách trình bày số liệu kiểm thử giữa thành viên.


Hướng cải thiện là chốt rõ phạm vi và dữ liệu dùng chung trước khi sửa, kiểm tra bản đóng gói trước khi tích hợp, lưu minh chứng có ngày và phiên bản, đồng thời kiểm tra lại các lỗi đã sửa. Các nội dung này đã được trình bày trong buổi Sprint Review ngày 03/11/2026. Nhóm lưu ảnh trình diễn và bản ghi buổi họp cùng kết quả tổng kết trên GitHub.


\begingroup
\small
\setlength{\tabcolsep}{5pt}
\renewcommand{\arraystretch}{1.2}
\begin{longtable}{|>{\raggedright\arraybackslash}p{\dimexpr 0.35932\linewidth-2\tabcolsep-2\arrayrulewidth\relax}|>{\raggedright\arraybackslash}p{\dimexpr 0.64068\linewidth-2\tabcolsep-2\arrayrulewidth\relax}|}
\hline
\textbf{Nội dung} & \textbf{Ghi nhận} \\ \hline
\endfirsthead
\hline
\textbf{Nội dung} & \textbf{Ghi nhận} \\ \hline
\endhead
\endfoot
\endlastfoot
Điều làm tốt & Chia việc theo từng task; tích hợp thành luồng mua hàng hoàn chỉnh; review mã nguồn và kiểm tra lại lỗi sau khi sửa; lưu ảnh và log minh chứng. \\ \hline
Khó khăn & Phối hợp giao diện với API; giữ thông tin đơn sau thanh toán; xử lý giao diện khi chạy Docker; bảo đảm kho và điểm không bị cập nhật lặp. \\ \hline
Cải thiện & Thống nhất dữ liệu và phạm vi trước khi sửa; trao đổi khi thay đổi tệp dùng chung; kiểm tra bản Docker trước khi tích hợp; ghi rõ phiên bản và môi trường kiểm thử. \\ \hline
Ngày họp và tham dự & Ngày 03/11/2026, qua Google Meet, ảnh ghi nhận 6 người tham gia. Hiển trình diễn; Hiệp đóng vai người nhận bàn giao. \\ \hline
Kết quả và phản hồi & Đã trình diễn luồng mua hàng, thanh toán lỗi rồi thử lại, thanh toán thành công và theo dõi đơn tại quầy. Buổi họp không ghi nhận góp ý bổ sung hoặc đầu việc mới. \\ \hline
\end{longtable}
\endgroup


\chapter{KẾT LUẬN VÀ BÀN GIAO SẢN PHẨM}

\Needspace{8\baselineskip}
\section{Kết quả đạt được}

BrewLite đã có bản tích hợp T01–T10, sử dụng đúng hai nền tảng Next.js và NestJS, lưu dữ liệu PostgreSQL và chạy bằng Docker Compose. Sản phẩm thực hiện được các bước chọn món, quản lý giỏ, xác thực, đặt đơn, thanh toán mô phỏng, xác nhận và lịch sử.


Nhóm đã tổng hợp mã nguồn, hướng dẫn chạy, lịch sử review và minh chứng kiểm thử. Buổi Sprint Review ngày 03/11/2026 đã trình diễn bản tích hợp và giới thiệu các tài liệu bàn giao, gồm kho GitHub, cấu hình Docker Compose và README.


\Needspace{8\baselineskip}
\section{Nội dung trình diễn và tiếp nhận phản hồi}

Ngày 03/11/2026, nhóm tổ chức Sprint Review qua Google Meet với 6 người tham gia được ghi nhận trên ảnh. Hiển chia sẻ màn hình và thao tác trên hệ thống; Hiệp đóng vai người nhận bàn giao. Nội dung gồm trình diễn sản phẩm, giới thiệu tài liệu chạy hệ thống và tổng kết quá trình làm việc.


\begin{enumerate}
    \item Trình diễn Menu, xem chi tiết đồ uống, chọn kích cỡ và số phần topping, thêm món vào giỏ và đăng nhập tài khoản.
    \item Kiểm tra giỏ hàng và số tiền phải trả; mô phỏng thanh toán thất bại, giữ giỏ hàng và quay lại thử thanh toán.
    \item Thanh toán thành công đơn số 26, gồm một ly cà phê sữa kích cỡ M với trân châu và kem sữa, tổng cộng 52.000đ. Giao diện ghi nhận điểm tăng từ 11 lên 16 và giỏ hàng về 0.
    \item Mở trang nhân viên để theo dõi đơn; xem lại chi tiết đơn với trạng thái \texttt{COMPLETED}.
    \item Giới thiệu kho mã nguồn GitHub, README và bộ chạy Docker Compose; trình bày việc làm tốt, khó khăn, hạn chế và cách cải thiện.
\end{enumerate}


\textbf{Kết quả:} Các chức năng trên đã được trình diễn trong buổi họp; luồng thanh toán thất bại có thể thử lại và hoàn tất đơn hàng thành công.


\textbf{Phản hồi:} Buổi họp không ghi nhận góp ý bổ sung hoặc đầu việc mới.


\textbf{Minh chứng:} Nhóm lưu 10 ảnh minh họa nội dung trình diễn và bản ghi buổi họp khoảng 14 phút 31 giây. Do phần hình của bản ghi gặp lỗi, ảnh được sử dụng để minh họa thao tác. Kết quả tổng kết và liên kết minh chứng được đăng tại Issue \#11.


\Needspace{8\baselineskip}
\section{Hạn chế và hướng phát triển}

Bản A1 dùng cổng thanh toán mô phỏng, chưa kết nối dịch vụ thanh toán thật. Kết quả hiệu năng hiện chỉ thuộc môi trường cục bộ và dữ liệu mẫu. Biểu đồ Burndown được tái dựng theo số task còn lại từ lịch sử Issue, chưa có dữ liệu điểm ước lượng hằng ngày.


Nếu tiếp tục phát triển, nhóm có thể khảo sát thanh toán thật, mở rộng theo dõi vận hành và kiểm thử tải. Đây là hướng đề xuất cho phiên bản sau, chưa phải phần đã hoàn thành hoặc cam kết trong A1.


\Needspace{8\baselineskip}
\section{Tài liệu đối chiếu và minh chứng}

Tài liệu chính gồm đề BrewLite Công nghệ Phần mềm phiên bản 1.0, mẫu báo cáo SGU, nhật ký task và lịch sử mã nguồn. Minh chứng kiểm thử là ảnh, log và báo cáo do từng thành viên gửi trên Issue \#11; người đọc có thể đối chiếu theo ngày và bản mã nguồn.


Đề bài và biểu mẫu: BrewLite Công nghệ Phần mềm VER 1.0, học kỳ I năm học 2026–2027; mẫu báo cáo SGU và ảnh bảng nhật ký giảng viên hướng dẫn.


Kho mã nguồn và lịch sử phát triển: \url{https://github.com/mhien06/BrewLite_A1}


Kết quả kiểm thử, Sprint Review và minh chứng bàn giao của nhóm: \url{https://github.com/mhien06/BrewLite_A1/issues/11}


\end{document}